\documentclass[11pt,a4paper]{article}

\usepackage[utf8]{inputenc}
\usepackage[italian]{babel}
\usepackage{amsmath, amssymb, amsthm}
\usepackage{graphicx}
\usepackage{xcolor}
\usepackage{tcolorbox}
\usepackage{geometry}
\geometry{a4paper, margin=2.5cm}

% Configurazione per i riquadri dei teoremi e definizioni
\tcbuselibrary{theorems}

\newtcbtheorem[number within=section]{teorema}{Teorema}%
{colback=blue!5,colframe=blue!50!black,fonttitle=\bfseries}{th}

\newtcbtheorem[number within=section]{definizione}{Definizione}%
{colback=green!5,colframe=green!50!black,fonttitle=\bfseries}{def}

\newtcolorbox{intuizione}{
    colback=yellow!10,
    colframe=orange!80!black,
    title=\textbf{Intuizione e Significato Pratico}
}

\title{\textbf{Riassunto Completo per l'Esame Orale}\\
\Large Probabilità Discreta e Processi Stocastici}
\author{Sintesi Estesa dalla Dispensa del Prof. S. Federico}
\date{Settembre 2026}

\begin{document}

\maketitle
\begin{abstract}
Questo documento costituisce un riassunto esteso e completo dell'intero programma d'esame. Include le definizioni matematiche rigorose, i teoremi fondamentali (inclusi argomenti avanzati come passeggiate aleatorie, martingala, 0-1 di Kolmogorov, e finanza) e le relative dimostrazioni, epurati del testo discorsivo.
\end{abstract}
\tableofcontents
\newpage

% ==========================================
% CAPITOLO 1
% ==========================================
\section{Teoria della Probabilità in Spazi Finiti}

\begin{definizione}{Spazio di Probabilità}{}
Uno spazio di probabilità finito è una terna $(\Omega, \mathcal{F}, \mathbb{P})$, dove:
\begin{itemize}
    \item $\Omega$ (Sample Space) è un insieme finito di esiti elementari.
    \item $\mathcal{F}$ è l'insieme degli eventi (sottoinsiemi di $\Omega$, spesso $\mathcal{F} = 2^\Omega$).
    \item $\mathbb{P}: \mathcal{F} \to [0,1]$ è una misura di probabilità tale che $\mathbb{P}(\Omega) = 1$ e per ogni $A, B \in \mathcal{F}$ disgiunti, $\mathbb{P}(A \cup B) = \mathbb{P}(A) + \mathbb{P}(B)$.
\end{itemize}
\end{definizione}

\begin{teorema}{Legge delle Probabilità Totali e Teorema di Bayes}{}
Sia $\{B_1, \dots, B_n\}$ una partizione di $\Omega$ (eventi disgiunti la cui unione è $\Omega$) con $\mathbb{P}(B_i) > 0$. Per ogni evento $A$:
\begin{enumerate}
    \item \textbf{Probabilità Totale:} $\mathbb{P}(A) = \sum_{i=1}^n \mathbb{P}(A \mid B_i)\mathbb{P}(B_i)$
    \item \textbf{Teorema di Bayes:} $\mathbb{P}(B_k \mid A) = \frac{\mathbb{P}(A \mid B_k)\mathbb{P}(B_k)}{\sum_{i=1}^n \mathbb{P}(A \mid B_i)\mathbb{P}(B_i)}$
\end{enumerate}
\end{teorema}
\paragraph{Dimostrazione (Bayes):} Dalla definizione di probabilità condizionata $\mathbb{P}(B_k \mid A) = \frac{\mathbb{P}(A \cap B_k)}{\mathbb{P}(A)}$. Poiché $\mathbb{P}(A \cap B_k) = \mathbb{P}(A \mid B_k)\mathbb{P}(B_k)$ e il denominatore si espande usando la legge della probabilità totale, segue la tesi. \qed

\subsection{Variabili Aleatorie e Valore Atteso}

\begin{definizione}{Valore Atteso e Varianza}{}
Sia $X: \Omega \to \mathbb{R}$ una Variabile Aleatoria (V.A.).
\begin{itemize}
    \item \textbf{Valore Atteso:} $\mathbb{E}[X] = \sum_{x} x \mathbb{P}(X=x)$.
    \item \textbf{Varianza:} $\text{Var}(X) = \mathbb{E}[(X - \mathbb{E}[X])^2] = \mathbb{E}[X^2] - (\mathbb{E}[X])^2$.
\end{itemize}
\end{definizione}

\begin{teorema}{Disuguaglianze di Markov e Chebyshev}{}
Sia $X$ una V.A.
\begin{enumerate}
    \item \textbf{Markov:} Se $X \ge 0$, per ogni $a > 0$, $\mathbb{P}(X \ge a) \le \frac{\mathbb{E}[X]}{a}$.
    \item \textbf{Chebyshev:} Per ogni $\epsilon > 0$, $\mathbb{P}(|X - \mathbb{E}[X]| \ge \epsilon) \le \frac{\text{Var}(X)}{\epsilon^2}$.
\end{enumerate}
\end{teorema}
\paragraph{Dimostrazione (Markov):} Osserviamo che $a \cdot \mathbf{1}_{\{X \ge a\}} \le X$. Prendendo il valore atteso da entrambe le parti si ottiene $a \mathbb{P}(X \ge a) \le \mathbb{E}[X]$, da cui la tesi. \qed

\newpage
% ==========================================
% CAPITOLO 2
% ==========================================
\section{Il Caso Discreto: da Finito a Numerabile}

Quando l'insieme degli esiti $\Omega$ è infinito numerabile, le definizioni si estendono tramite le serie. 

\begin{definizione}{Distribuzioni Fondamentali su base numerabile}{}
\begin{itemize}
    \item \textbf{Poisson($\lambda$):} $\mathbb{P}(X=k) = e^{-\lambda} \frac{\lambda^k}{k!}$ (per $k \in \mathbb{N}$). Modella eventi rari. $\mathbb{E}[X] = \lambda$, $\text{Var}(X) = \lambda$.
    \item \textbf{Geometrica($p$):} $\mathbb{P}(X=k) = (1-p)^{k-1} p$ (per $k \in \mathbb{N}\setminus\{0\}$). Modella il tempo del primo successo. $\mathbb{E}[X] = 1/p$. Gode della proprietà di \textbf{assenza di memoria}: $\mathbb{P}(X > n+m \mid X > n) = \mathbb{P}(X > m)$.
\end{itemize}
\end{definizione}

\subsection{Convergenza di Variabili Aleatorie}
Consideriamo una successione di V.A. $\{X_n\}$.
\begin{itemize}
    \item \textbf{Quasi Certamente (a.s.):} $\mathbb{P}(\{\omega : \lim_{n \to \infty} X_n(\omega) = X(\omega)\}) = 1$.
    \item \textbf{In Probabilità:} Per ogni $\epsilon > 0$, $\lim_{n \to \infty} \mathbb{P}(|X_n - X| > \epsilon) = 0$.
    \item \textbf{In Media (o in $L^1$):} $\lim_{n \to \infty} \mathbb{E}[|X_n - X|] = 0$.
    \item \textbf{In Legge (o Distribuzione):} Le funzioni di ripartizione convergono nei punti di continuità.
\end{itemize}

\begin{teorema}{Relazioni tra Convergenze}{}
Valgono le seguenti implicazioni (non invertibili in generale):
\[ \text{Quasi certamente} \implies \text{In Probabilità} \implies \text{In Legge} \]
\[ \text{In Media } L^1 \implies \text{In Probabilità} \]
\end{teorema}

\begin{teorema}{Legge Forte dei Grandi Numeri (SLLN)}{}
Sia $\{X_n\}_{n \ge 1}$ una successione di V.A. i.i.d. con $\mathbb{E}[|X_1|] < \infty$ e $\mathbb{E}[X_1] = \mu$. Detta $S_n = \sum_{i=1}^n X_i$, si ha:
\[ \frac{S_n}{n} \xrightarrow{a.s.} \mu \]
\end{teorema}
\begin{intuizione}
La media empirica di un campione converge "con certezza assoluta" al valore atteso teorico quando la dimensione del campione tende all'infinito. Questo è il teorema fondamentale della statistica e delle simulazioni Monte Carlo.
\end{intuizione}

\newpage
% ==========================================
% CAPITOLO 3
% ==========================================
\section{Processi Stocastici a Tempo Discreto}

\begin{definizione}{Filtrazione e Processo Adattato}{}
Una \textbf{filtrazione} $\mathbb{F} = \{\mathcal{F}_n\}_{n \ge 0}$ è una famiglia crescente di $\sigma$-algebre ($\mathcal{F}_0 \subseteq \mathcal{F}_1 \subseteq \dots$). Rappresenta il flusso di informazione col passare del tempo.
Un processo stocastico $\{X_n\}$ si dice \textbf{adattato} alla filtrazione $\mathbb{F}$ se, per ogni $n$, $X_n$ è $\mathcal{F}_n$-misurabile (ovvero il valore di $X_n$ è perfettamente noto avendo l'informazione al tempo $n$).
\end{definizione}

\begin{definizione}{Tempo di Arresto (Stopping Time)}{}
Una V.A. $\tau: \Omega \to \mathbb{N} \cup \{\infty\}$ è un tempo di arresto rispetto a $\mathbb{F}$ se per ogni $n \ge 0$, l'evento $\{\tau \le n\} \in \mathcal{F}_n$.
\end{definizione}
\begin{intuizione}
Un tempo di arresto è una regola per fermarsi che non richiede la chiaroveggenza. Puoi decidere di fermarti al tempo $n$ guardando solo ciò che è successo fino al tempo $n$. "Mi fermo alla prima volta che ottengo testa" è un tempo di arresto. "Mi fermo al lancio prima di ottenere testa" NON lo è.
\end{intuizione}

\begin{teorema}{Identità di Wald}{}
Sia $\{X_n\}_{n \ge 1}$ una sequenza di V.A. i.i.d. con $\mathbb{E}[|X_1|] < \infty$. Sia $N$ un tempo di arresto integrabile rispetto alla filtrazione generata dalle $X_n$, e supponiamo che $\mathbb{E}[N] < \infty$. Allora, posto $S_N = \sum_{n=1}^N X_n$:
\[ \mathbb{E}[S_N] = \mathbb{E}[N] \cdot \mathbb{E}[X_1] \]
\end{teorema}
\paragraph{Dimostrazione:}
Possiamo scrivere $S_N = \sum_{n=1}^\infty X_n \mathbf{1}_{\{N \ge n\}}$. Osserviamo che l'evento $\{N \ge n\} = \{N \le n-1\}^c \in \mathcal{F}_{n-1}$. Poiché $X_n$ è indipendente da $\mathcal{F}_{n-1}$ (le prove sono i.i.d.), $\mathbb{E}[X_n \mathbf{1}_{\{N \ge n\}}] = \mathbb{E}[X_n]\mathbb{E}[\mathbf{1}_{\{N \ge n\}}]$.
Sommando su $n$:
\[ \mathbb{E}[S_N] = \sum_{n=1}^\infty \mathbb{E}[X_1]\mathbb{P}(N \ge n) = \mathbb{E}[X_1] \sum_{n=1}^\infty \mathbb{P}(N \ge n) = \mathbb{E}[X_1]\mathbb{E}[N] \]
\qed

\newpage
% ==========================================
% CAPITOLO 4
% ==========================================
\section{Martingale}

\begin{definizione}{Martingala, Sub/Super-martingala}{}
Sia $\{M_n\}_{n \ge 0}$ un processo adattato a $\{\mathcal{F}_n\}$ con $\mathbb{E}[|M_n|] < \infty$.
\begin{itemize}
    \item \textbf{Martingala:} $\mathbb{E}[M_{n+1} \mid \mathcal{F}_n] = M_n$ (gioco equo).
    \item \textbf{Submartingala:} $\mathbb{E}[M_{n+1} \mid \mathcal{F}_n] \ge M_n$ (trend crescente, gioco favorevole).
    \item \textbf{Supermartingala:} $\mathbb{E}[M_{n+1} \mid \mathcal{F}_n] \le M_n$ (trend decrescente, gioco a perdere).
\end{itemize}
\end{definizione}

\begin{teorema}{Teorema di Arresto Opzionale (Optional Stopping Theorem di Doob)}{}
Sia $\{M_n\}$ una martingala e $\tau$ un tempo di arresto. Sotto opportune condizioni di regolarità (ad esempio, $\tau$ è limitato quasi certamente, $\tau \le K$), si ha:
\[ \mathbb{E}[M_\tau] = \mathbb{E}[M_0] \]
\end{teorema}
\paragraph{Dimostrazione (per $\tau$ limitato):}
Definiamo il processo arrestato $M_{n \wedge \tau} = M_n \mathbf{1}_{\{n \le \tau\}} + M_\tau \mathbf{1}_{\{n > \tau\}}$.
Si può facilmente dimostrare che se $M_n$ è una martingala, anche il processo arrestato $M_{n \wedge \tau}$ è una martingala rispetto a $\mathcal{F}_n$.
Poiché è una martingala, ha valore atteso costante: $\mathbb{E}[M_{n \wedge \tau}] = \mathbb{E}[M_0]$ per ogni $n$.
Poiché $\tau$ è limitato da una costante deterministica $K$, prendendo $n = K$ otteniamo:
$M_{K \wedge \tau} = M_\tau$. Quindi $\mathbb{E}[M_\tau] = \mathbb{E}[M_0]$.
\qed

\begin{intuizione}
In un casinò che offre solo scommesse eque, nessuna strategia (tempo di arresto) basata sull'osservazione del passato può permettere a un giocatore di avere un guadagno atteso positivo. Non si può creare valore dal nulla fermandosi "al momento giusto".
\end{intuizione}

\newpage
% ==========================================
% CAPITOLO 5
% ==========================================
\section{Passeggiate Aleatorie (Random Walks)}

\begin{teorema}{Lemmi di Borel-Cantelli}{}
Sia $\{A_n\}$ una successione di eventi in uno spazio di probabilità. Consideriamo l'evento $\limsup A_n$ (cioè l'evento che accadano infiniti eventi $A_n$).
\begin{enumerate}
    \item \textbf{Primo Lemma:} Se $\sum_{n=1}^\infty \mathbb{P}(A_n) < \infty$, allora $\mathbb{P}(\limsup A_n) = 0$.
    \item \textbf{Secondo Lemma:} Se $\sum_{n=1}^\infty \mathbb{P}(A_n) = \infty$ e gli eventi $A_n$ sono \textbf{indipendenti}, allora $\mathbb{P}(\limsup A_n) = 1$.
\end{enumerate}
\end{teorema}

\begin{teorema}{Legge 0-1 di Kolmogorov}{}
Sia $\{X_n\}$ una successione di V.A. indipendenti. Definiamo la $\sigma$-algebra "coda" $\mathcal{T} = \bigcap_{n=1}^\infty \sigma(X_n, X_{n+1}, \dots)$. Per ogni evento asintotico $A \in \mathcal{T}$, la probabilità $\mathbb{P}(A)$ è esattamente $0$ oppure $1$.
\end{teorema}

\subsection{Rovina del Giocatore (Gambler's Ruin)}
Sia $S_n = \sum_{i=1}^n X_i$ una passeggiata aleatoria su $\mathbb{Z}$ con $X_i \in \{-1, +1\}$ e $\mathbb{P}(X_i = 1) = p$.
Un giocatore inizia con un capitale $x$ e l'obiettivo è raggiungere il target $N$ prima di finire in rovina a $0$.
\begin{teorema}{Probabilità di Vittoria}{}
Sia $P_x$ la probabilità di raggiungere $N$ partendo da $x$.
\[
P_x = 
\begin{cases}
\frac{x}{N} & \text{se il gioco è equo } (p = 1/2) \\
\frac{1 - ((1-p)/p)^x}{1 - ((1-p)/p)^N} & \text{se } p \neq 1/2
\end{cases}
\]
\end{teorema}

\paragraph{Dimostrazione (tramite Martingale, caso $p=1/2$):}
Se il gioco è equo, $\mathbb{E}[X_i] = 0$. Di conseguenza la passeggiata aleatoria $S_n$ è una martingala (infatti $\mathbb{E}[S_{n+1} \mid \mathcal{F}_n] = S_n + \mathbb{E}[X_{n+1}] = S_n$).
Definiamo il tempo di arresto $\tau = \min\{n \ge 0 : S_n = 0 \text{ o } S_n = N\}$. È noto che $\tau < \infty$ quasi certamente in una passeggiata aleatoria in dominio limitato.
Essendo $S_{n \wedge \tau}$ una martingala limitata nell'intervallo $[0, N]$, il Teorema di Arresto Opzionale garantisce:
\[ x = \mathbb{E}[S_0] = \mathbb{E}[S_\tau] \]
Sappiamo che al tempo $\tau$, il processo vale $N$ (con probabilità $P_x$) oppure $0$ (con probabilità $1-P_x$).
Dunque:
\[ \mathbb{E}[S_\tau] = N \cdot P_x + 0 \cdot (1-P_x) = N \cdot P_x \]
Uguagliando, $x = N P_x$, da cui $P_x = x/N$.
\qed

\newpage
% ==========================================
% CAPITOLO 6
% ==========================================
\section{Modelli Finanziari in Spazi Finiti}

\begin{definizione}{Arbitraggio}{}
Un portafoglio autofinanziante (senza iniezione o prelievo di fondi) che parte da un valore iniziale nullo $V_0 = 0$, ma garantisce un valore finale quasi certamente non negativo ($\mathbb{P}(V_T \ge 0) = 1$) e positivo con probabilità non nulla ($\mathbb{P}(V_T > 0) > 0$), costituisce una \textbf{opportunità di arbitraggio} ("pasto gratis").
\end{definizione}

\begin{teorema}{Primo Teorema Fondamentale del Pricing (FTAP)}{}
Un modello di mercato finanziario a tempo discreto in uno spazio finito è privo di opportunità di arbitraggio se e solo se esiste almeno una \textbf{Misura Martingala Equivalente} (EMM) $\mathbb{Q}$.
Sotto $\mathbb{Q}$, i prezzi scontati degli asset sono martingale: 
\[ \mathbb{E}^{\mathbb{Q}}[ \tilde{S}_{n+1} \mid \mathcal{F}_n ] = \tilde{S}_n \]
\end{teorema}

\begin{teorema}{Secondo Teorema Fondamentale (Completezza)}{}
Un mercato privo di arbitraggi è \textbf{completo} (ovvero ogni derivato finanziario può essere replicato perfettamente tramite un portafoglio) se e solo se la misura martingala equivalente $\mathbb{Q}$ è \textbf{unica}.
\end{teorema}

\subsection{Modello Binomiale di Cox-Ross-Rubinstein (CRR)}
Nel modello binomiale uni-periodale, lo stock al tempo 0 vale $S_0$. Al tempo 1 può salire a $S_0(1+u)$ o scendere a $S_0(1+d)$, con $d < r < u$ (dove $r$ è il tasso privo di rischio). 
\begin{teorema}{Prezzaggio nel Modello Binomiale}{}
L'unica probabilità neutrale al rischio $\mathbb{Q}$ assegna probabilità $q$ al rialzo, dove:
\[ q = \frac{r - d}{u - d} \]
Il prezzo equo di un derivato con payoff $C_u$ (se lo stock sale) e $C_d$ (se scende) è l'aspettativa attualizzata sotto la misura $\mathbb{Q}$:
\[ C_0 = \frac{1}{1+r} [ q C_u + (1-q) C_d ] \]
\end{teorema}

\newpage
% ==========================================
% CAPITOLO 7
% ==========================================
\section{Catene di Markov (Nucleo del Progetto SIR)}

\begin{definizione}{Catena di Markov e Proprietà (Strong)}{}
Un processo stocastico $\{X_n\}$ a valori in uno spazio discreto $S$ è una \textbf{Catena di Markov} se il passato e il futuro sono condizionatamente indipendenti dato il presente:
\[ \mathbb{P}(X_{n+1} = j \mid X_n = i, X_{n-1} = i_{n-1}, \dots) = \mathbb{P}(X_{n+1} = j \mid X_n = i) = p_{ij} \]
Il processo soddisfa la \textbf{Proprietà di Markov Forte} se la proprietà di assenza di memoria vale anche rimpiazzando il tempo deterministico $n$ con un qualsiasi \textit{tempo di arresto} quasi certamente finito.
\end{definizione}

\begin{teorema}{Equazioni di Chapman-Kolmogorov}{}
Le probabilità di transizione in $n+m$ passi soddisfano:
\[ p_{ij}^{(n+m)} = \sum_{k \in S} p_{ik}^{(n)} p_{kj}^{(m)} \quad \iff \quad P^{(n+m)} = P^{(n)} P^{(m)} \]
\end{teorema}
\paragraph{Dimostrazione:} Disgiungendo lo spazio su tutti i possibili stati intermedio $k$ assunti al tempo $n$ e applicando la legge della probabilità totale e l'assenza di memoria di Markov. \qed

\subsection{Distribuzioni Invarianti ed Ergodicità}

Una distribuzione (vettore riga) $\pi$ è \textbf{invariante} se $\pi P = \pi$.
Una Catena è \textbf{Irriducibile} se esiste un percorso con probabilità positiva tra ogni coppia di stati. È \textbf{Aperiodica} se il massimo comun divisore delle lunghezze dei percorsi chiusi (loop) di ogni stato è 1.

\begin{teorema}{Teorema Ergodico per Catene Finite}{}
Sia $\{X_n\}$ una Catena di Markov a stati finiti, irriducibile e aperiodica. Allora:
\begin{enumerate}
    \item Esiste un'unica distribuzione invariante $\pi$ strettamente positiva ($\pi_i > 0$).
    \item Le probabilità di transizione convergono asintoticamente alla misura invariante, indipendentemente dallo stato iniziale:
    \[ \lim_{n \to \infty} p_{ij}^{(n)} = \pi_j \]
    \item La distribuzione invariante è inversamente proporzionale al tempo medio di ritorno nello stato $j$ (chiamato $m_j$): $\pi_j = \frac{1}{m_j}$.
\end{enumerate}
\end{teorema}
\begin{intuizione}
Questa è la base matematica del "ciclo della vita" della Catena: in un sistema chiuso e connesso (irriducibile), a lungo termine si instaura un equilibrio dinamico perfetto, dove la probabilità di trovarsi in un certo stato rappresenta esattamente la percentuale di tempo che il sistema passa in quello stato.
\end{intuizione}

\newpage
% ==========================================
% CAPITOLO 8
% ==========================================
\section{Controllo Ottimo Stocastico a Tempo Discreto}

Nel Controllo Ottimo, l'agente osserva lo stato $x_n$ del sistema e sceglie un'azione (controllo) $u_n$ per massimizzare (o minimizzare) una funzione obiettivo attesa.
L'evoluzione del sistema è descritta dalla dinamica:
\[ X_{n+1} = f(X_n, u_n, W_n) \]
dove $W_n$ è un rumore aleatorio (disturbo) indipendente.

L'obiettivo (su orizzonte infinito con fattore di sconto $\beta \in (0,1)$) è trovare una \textit{policy} di controllo (una funzione $u_n = g(X_n)$) che massimizzi:
\[ V(x) = \mathbb{E} \left[ \sum_{n=0}^\infty \beta^n L(X_n, u_n, W_n) \bigg| X_0 = x \right] \]
dove $L$ è il ricavo/costo istantaneo e $V(x)$ è la \textbf{Funzione Valore} ottima partendo da $x$.

\begin{teorema}{Equazione di Ottimalità di Bellman}{}
La funzione valore ottima $V(x)$ per il problema stazionario a orizzonte infinito deve soddisfare l'equazione ricorsiva di Bellman:
\[ V(x) = \max_u \mathbb{E} \left[ L(x, u, W_0) + \beta V(f(x, u, W_0)) \right] \]
\end{teorema}
\paragraph{Dimostrazione (Idea del Principio di Ottimalità):}
Il ricavo totale atteso a partire dal tempo $0$ può essere spezzato nella somma del ricavo immediato al tempo $0$ e nel ricavo atteso dal tempo $1$ in poi (scontato di un fattore $\beta$). Se stiamo seguendo una policy ottima, dal tempo $1$ in poi (avendo raggiunto il nuovo stato stocastico $f(x, u, W_0)$) dovremo comportarci in modo ottimale per forza, da cui compare $V(f(x, u, W_0))$. Prendendo il massimo sulla scelta iniziale $u$, si ottiene l'equazione. \qed

\begin{intuizione}
\textbf{Principio di Ottimalità di Bellman:} Indipendentemente dallo stato iniziale e dalle prime decisioni prese, le decisioni rimanenti devono necessariamente costituire una policy ottima rispetto allo stato derivante dalla primissima decisione. 
Invece di calcolare tutte le possibili sequenze di mosse infinite, l'equazione di Bellman trasforma il problema nel trovare la "mossa successiva migliore" assumendo di comportarsi perfettamente da domani in poi.
\end{intuizione}

\end{document}
